\section{Time Hierarchy --- Clock Machine}
%%% generated by the blueprint pipeline from TCSlib.Complexity.TimeHierarchy.ClockMachine
\subsection*{Time Hierarchy --- Clock Machine: Introduction}

Declarations of the Lean module \texttt{TCSlib.Complexity.TimeHierarchy.ClockMachine}.

\subsection*{Time Hierarchy --- Clock Machine: Results}

\begin{definition}[Output register states]
\lean{Complexity.TimeHierarchy.OutReg}
\label{Complexity.TimeHierarchy.OutReg}
\leanok
The output register has three possible forms: empty, containing one Boolean symbol, or
containing multiple symbols.
\end{definition}

\begin{definition}[Register state from an output string]
\lean{Complexity.TimeHierarchy.OutReg.ofList}
\label{Complexity.TimeHierarchy.OutReg.ofList}
\leanok
\uses{Complexity.TimeHierarchy.OutReg}
For a finite Boolean string, its register state is empty if the string has length zero,
records its symbol if the string has length one, and is the multiple-symbol state if its
length is at least two.
\end{definition}

\begin{definition}[Appending to an output register]
\lean{Complexity.TimeHierarchy.OutReg.push}
\label{Complexity.TimeHierarchy.OutReg.push}
\leanok
\uses{Complexity.TimeHierarchy.OutReg}
Appending no symbol leaves the register state unchanged. Appending a Boolean symbol to
an empty register records that symbol; appending one to a register containing one symbol
changes the state to multiple symbols; and appending one to a register already in the
multiple-symbol state leaves it there.
\end{definition}

\begin{lemma}[Output register under appending a symbol]
\lean{Complexity.TimeHierarchy.OutReg.ofList_push}
\label{Complexity.TimeHierarchy.OutReg.ofList_push}
\leanok
\uses{Complexity.TimeHierarchy.OutReg.ofList, Complexity.TimeHierarchy.OutReg.push}
\difficulty{4}
For every finite Boolean word $w$ and optional Boolean symbol $o$, the output register
of the word formed by appending $o$ to $w$ equals the result of pushing $o$ onto the
output register of $w$.
\end{lemma}

\begin{lemma}[Recognition of the singleton true output]
\lean{Complexity.TimeHierarchy.OutReg.ofList_eq_one_true}
\label{Complexity.TimeHierarchy.OutReg.ofList_eq_one_true}
\leanok
\uses{Complexity.TimeHierarchy.OutReg.ofList}
\difficulty{4}
For every finite Boolean word $w$, its output register is the singleton register for
true exactly when $w$ is the one-symbol word containing true.
\end{lemma}

\begin{definition}[Binary value of a Boolean list]
\lean{Complexity.TimeHierarchy.ctrVal}
\label{Complexity.TimeHierarchy.ctrVal}
\leanok
For a list of Booleans, define its value recursively by assigning the empty list value
$0$ and assigning a list with first bit $b$ and remaining list $bs$ the value $(1	ext{
if }b	ext{ is true, otherwise }0)+2$ times the value of $bs$. Thus the first list
element is the least significant bit.
\end{definition}

\begin{definition}[Number of true bits]
\lean{Complexity.TimeHierarchy.ctrPop}
\label{Complexity.TimeHierarchy.ctrPop}
\leanok
For a list of Booleans, define its population count recursively: the empty list has
count $0$, and a list with first bit $b$ and remaining list $bs$ has count $(1	ext{ if
}b	ext{ is true, otherwise }0)$ plus the count of $bs$.
\end{definition}

\begin{lemma}[True Bits Bounded by Length]
\lean{Complexity.TimeHierarchy.ctrPop_le_length}
\label{Complexity.TimeHierarchy.ctrPop_le_length}
\leanok
\uses{Complexity.TimeHierarchy.ctrPop}
\difficulty{4}
For every finite Boolean list $s$, the number of entries of $s$ equal to true is at most
the length of $s$.
\end{lemma}

\begin{lemma}[Value of the Binary Expansion]
\lean{Complexity.TimeHierarchy.ctrVal_bits}
\label{Complexity.TimeHierarchy.ctrVal_bits}
\leanok
\uses{Complexity.TimeHierarchy.ctrVal}
\difficulty{4}
For every natural number $n$, interpreting its binary digits with the first digit as the
least significant digit gives $n$.
\end{lemma}

\begin{lemma}[Value After False Digits]
\lean{Complexity.TimeHierarchy.ctrVal_replicate_false}
\label{Complexity.TimeHierarchy.ctrVal_replicate_false}
\leanok
\uses{Complexity.TimeHierarchy.ctrVal}
\difficulty{4}
For every natural number $j$ and Boolean list $r$, the value of the list formed by
placing $j$ false digits before $r$ is $2^j$ times the value of $r$.
\end{lemma}

\begin{lemma}[Value After True Digits]
\lean{Complexity.TimeHierarchy.ctrVal_replicate_true}
\label{Complexity.TimeHierarchy.ctrVal_replicate_true}
\leanok
\uses{Complexity.TimeHierarchy.ctrVal}
\difficulty{4}
For every natural number $j$ and Boolean list $r$, the value of the list formed by
placing $j$ true digits before $r$, plus one, equals $2^j$ times the value of $r$ plus
one.
\end{lemma}

\begin{lemma}[Binary Borrow Identity]
\lean{Complexity.TimeHierarchy.ctrVal_borrow}
\label{Complexity.TimeHierarchy.ctrVal_borrow}
\leanok
\uses{Complexity.TimeHierarchy.ctrVal, Complexity.TimeHierarchy.ctrVal_replicate_false, Complexity.TimeHierarchy.ctrVal_replicate_true}
\difficulty{3}
For every natural number $j$ and Boolean list $r$, the value of the list consisting of
$j$ true digits, then a false digit, then $r$, plus one equals the value of the list
consisting of $j$ false digits, then a true digit, then $r$.
\end{lemma}

\begin{lemma}[True Bits in a Repeated Prefix]
\lean{Complexity.TimeHierarchy.ctrPop_replicate}
\label{Complexity.TimeHierarchy.ctrPop_replicate}
\leanok
\uses{Complexity.TimeHierarchy.ctrPop}
\difficulty{4}
For every natural number $j$, Boolean value $b$, and Boolean list $r$, the number of
true entries in the list formed by placing $j$ copies of $b$ before $r$ equals $j$ if
$b$ is true, and otherwise equals the number of true entries in $r$.
\end{lemma}

\begin{lemma}[First True Digit in a Positive-Value List]
\lean{Complexity.TimeHierarchy.ctrVal_pos_decomp}
\label{Complexity.TimeHierarchy.ctrVal_pos_decomp}
\leanok
\uses{Complexity.TimeHierarchy.ctrVal}
\difficulty{4}
For every Boolean list $s$ whose binary value is nonzero, there are a natural number $j$
and a Boolean list $r$ such that $s$ consists of $j$ false digits followed by a true
digit and then $r$.
\end{lemma}

\begin{lemma}[Zero Value Characterization]
\lean{Complexity.TimeHierarchy.ctrVal_zero_eq}
\label{Complexity.TimeHierarchy.ctrVal_zero_eq}
\leanok
\uses{Complexity.TimeHierarchy.ctrVal}
\difficulty{4}
For every Boolean list $s$ whose binary value is zero, $s$ consists entirely of false
digits; equivalently, it is the list of false values whose length is the length of $s$.
\end{lemma}

\begin{definition}[States of a Clocked Machine]
\lean{Complexity.TimeHierarchy.ClockState}
\label{Complexity.TimeHierarchy.ClockState}
\leanok
\uses{Complexity.TimeHierarchy.OutReg}
Let $K$ and $W$ be types, and let $R$ be a type of output-register values. A clock state
is one of six forms: a run state carrying an optional element of $K$; a counter-scan
state; an input-start state; an input-scan state; a decrement state carrying an element
of $W$ and an element of $R$; or a return state carrying an element of $W$ and an
element of $R$.
\end{definition}

\begin{definition}[Counter tape index]
\lean{Complexity.TimeHierarchy.ctrIdx}
\label{Complexity.TimeHierarchy.ctrIdx}
\leanok
For natural numbers $k_K$ and $k_W$, consider $k_K+1+k_W$ tape indices arranged in three
consecutive blocks of lengths $k_K$, $1$, and $k_W$. The counter index is the unique
index in the middle, namely the index at offset $k_K$ from the beginning.
\end{definition}

\begin{definition}[Counter tape action]
\lean{Complexity.TimeHierarchy.ctrAction}
\label{Complexity.TimeHierarchy.ctrAction}
\leanok
\uses{Turing.Action, Turing.FinTM.tapeBlocks}
For natural numbers $k_K,k_W$, an optional optional Boolean $w$, a head movement
$d\in\{-1,0,1\}$, and an optional successor state $q$, define an action on
$k_K+1+k_W$ work tapes over the Boolean alphabet: the input head stays in place, every
work tape except the middle tape is left unchanged and stationary, and the middle tape
receives write instruction $w$ and moves by $d$. The action emits no output and passes
to $q$, halting if $q$ is absent.
\end{definition}

\begin{definition}[Clocked simulation transition]
\lean{Complexity.TimeHierarchy.clockTr}
\label{Complexity.TimeHierarchy.clockTr}
\leanok
\uses{Complexity.TimeHierarchy.ClockState, Complexity.TimeHierarchy.OutReg.push, Complexity.TimeHierarchy.ctrAction, Complexity.TimeHierarchy.ctrIdx, Turing.Action, Turing.FinTM, Turing.FinTM.controlAction, Turing.FinTM.tapeBlocks}
Let $K$ and $W$ be finite Turing machines over the Boolean alphabet. Given a clock
state, an optional input symbol, and the symbols scanned on the $K.k+1+W.k$ work tapes,
define a transition as follows. In a simulation state, perform $K$’s transition on the
first $K.k$ tapes, place its optional output on the middle tape, and leave the final
$W.k$ tapes unchanged; if $K$ halts, enter the counter scan. The counter scan moves
right and starts the input scan when the middle cell is blank, and otherwise moves left
and continues scanning. The input scan moves left until the input cell is blank, then
moves right and starts a decrement state with $W$’s initial state and an empty output
register. In a decrement state, a blank middle cell halts and emits true; a false middle
symbol is changed to true and the middle head moves right; a true middle symbol triggers
$W$’s transition on the final $W.k$ tapes, changes the middle symbol to false, and moves
the middle head left. If $W$ halts on that step, the transition emits true exactly when
the updated output register does not contain precisely one true bit; if $W$ continues,
the transition enters a return state carrying its new state and the updated register.
The return state moves right and enters the corresponding decrement state when the
middle cell is blank, and otherwise moves left and continues returning. Tapes not
specified as being written are left unchanged, and their heads remain stationary.
\end{definition}

\begin{definition}[Clocked machine construction]
\lean{Complexity.TimeHierarchy.clockTM}
\label{Complexity.TimeHierarchy.clockTM}
\leanok
\uses{Complexity.TimeHierarchy.ClockState, Complexity.TimeHierarchy.clockTr, Turing.FinTM}
Given finite Turing machines $K$ and $W$ over the Boolean alphabet, define a finite
Turing machine with $K.k+1+W.k$ work tapes, state set
$	ext{ClockState}(K.	ext{State},W.	ext{State})$, initial state $	ext{.kRun}(	ext{some
}K.q_0)$, and transition function $	ext{clockTr}(K,W)$.
\end{definition}

\begin{lemma}[Clock machine step in a state]
\lean{Complexity.TimeHierarchy.clockTM_step_some}
\label{Complexity.TimeHierarchy.clockTM_step_some}
\leanok
\uses{Complexity.TimeHierarchy.ClockState, Complexity.TimeHierarchy.clockTM, Complexity.TimeHierarchy.clockTr, Turing.Action.apply, Turing.Cfg, Turing.Cfg.inputSymbol, Turing.Cfg.workTapeSymbols, Turing.MultiTapeTM.step}
\difficulty{3}
Let $K$ and $W$ be finite Boolean Turing machines, and let $c$ be a configuration of
their clock machine on Boolean input $x$. If the state of $c$ is a state $q$, then one
step from $c$ equals applying to $c$ the action prescribed by the clock transition
function for $q$, the symbol under the input head, and the symbols under the work-tape
heads.
\end{lemma}

\begin{definition}[Configuration embedding]
\lean{Complexity.TimeHierarchy.kCfg}
\label{Complexity.TimeHierarchy.kCfg}
\leanok
\uses{Complexity.TimeHierarchy.clockTM, Turing.Cfg, Turing.FinTM.bufferTape, Turing.FinTM.tapeBlocks}
Let $K$ and $W$ be finite Turing machines over the Boolean alphabet, let $x$ be a
Boolean input string, and let $c$ be a configuration of $K$ on input $x$. The
corresponding configuration of the clocked machine on $x$ has simulation state carrying
the state of $c$, the same input-head position, and empty output. Its first $K.k$ work
tapes and head positions are those of $c$; the next work tape contains the output string
of $c$ and its head is at position equal to the length of that string; the remaining
$W.k$ work tapes are blank and their heads are at position $0$.
\end{definition}

\begin{lemma}[Initial lifted configuration]
\lean{Complexity.TimeHierarchy.kCfg_init}
\label{Complexity.TimeHierarchy.kCfg_init}
\leanok
\uses{Complexity.TimeHierarchy.clockTM, Complexity.TimeHierarchy.kCfg, Turing.FinTM.tapeBlocks, Turing.MultiTapeTM.initCfg}
\difficulty{5}
Let $K$ and $W$ be finite Turing machines over the binary alphabet, and let $x$ be a
binary input. The initial configuration of the clock machine built from $K$ and $W$ on
$x$ is the lifted initial configuration of $K$ on $x$.
\end{lemma}

\begin{lemma}[One-step preservation of lifted configurations]
\lean{Complexity.TimeHierarchy.kCfg_step}
\label{Complexity.TimeHierarchy.kCfg_step}
\leanok
\uses{Complexity.TimeHierarchy.clockTM, Complexity.TimeHierarchy.clockTM_step_some, Complexity.TimeHierarchy.clockTr, Complexity.TimeHierarchy.kCfg, Turing.Action.apply, Turing.Cfg, Turing.Cfg.inputSymbol, Turing.Cfg.workTapeSymbols, Turing.FinTM.bufferTape_append, Turing.MultiTapeTM.step}
\difficulty{6}
Let $K$ and $W$ be finite Turing machines over the binary alphabet, let $x$ be a binary
input, and let $c$ be a configuration of $K$ on $x$ with a present state. One step of
the clock machine from the lifted configuration of $c$ is the lift of the configuration
reached by one step of $K$ from $c$.
\end{lemma}

\begin{lemma}[Lifted computation before halting]
\lean{Complexity.TimeHierarchy.kCfg_run}
\label{Complexity.TimeHierarchy.kCfg_run}
\leanok
\uses{Complexity.TimeHierarchy.clockTM, Complexity.TimeHierarchy.kCfg, Complexity.TimeHierarchy.kCfg_step, Turing.Cfg, Turing.MultiTapeTM.runFrom, Turing.MultiTapeTM.runFrom_succ_eq_step'}
\difficulty{4}
Let $K$ and $W$ be finite Turing machines over the binary alphabet, let $x$ be a binary
input, and let $c$ be a configuration of $K$ on $x$. If the computation from $c$ has a
present state at every time $s<t$, then running the clock machine from the lifted
configuration of $c$ for $t$ steps gives the lifted configuration reached by $K$ from
$c$ after $t$ steps.
\end{lemma}

\begin{definition}[Clock configuration assembly]
\lean{Complexity.TimeHierarchy.ctrCfg}
\label{Complexity.TimeHierarchy.ctrCfg}
\leanok
\uses{Complexity.TimeHierarchy.ClockState, Complexity.TimeHierarchy.clockTM, Turing.Cfg, Turing.FinTM.tapeBlocks}
Let $K$ and $W$ be finite Boolean Turing machines, let $x$ be a Boolean input, and let
$c$ be a configuration of $W$ on $x$. Given a clock state $s$, a buffer tape $\tau$, a
buffer-head position $z\in\mathbb Z$, work-tape contents $kt$, work-head positions $kh$,
and an output word $out$, form a configuration of the clock machine with state $s$,
input-head position inherited from $c$, and output $out$. Its work tapes, in order, are
$kt$, the buffer tape $\tau$, and the work tapes of $c$; their head positions are $kh$,
$z$, and the work-head positions of $c$.
\end{definition}

\begin{lemma}[Counter tape scanned symbol]
\lean{Complexity.TimeHierarchy.ctrCfg_ctr}
\label{Complexity.TimeHierarchy.ctrCfg_ctr}
\leanok
\uses{Complexity.TimeHierarchy.ClockState, Complexity.TimeHierarchy.ctrCfg, Complexity.TimeHierarchy.ctrIdx, Turing.Cfg, Turing.Cfg.workTapeSymbols}
\difficulty{1}
Let $K$ and $W$ be finite Boolean Turing machines, and form a clock configuration from
any clock state, any configuration of $W$, buffer-tape content $\tau$, buffer-head
position $z$, fixed contents and head positions for the $K$ work tapes, and output word.
The symbol scanned on the distinguished buffer tape is $\tau(z)$.
\end{lemma}

\begin{lemma}[Right work-tape symbols in assembled configuration]
\lean{Complexity.TimeHierarchy.ctrCfg_right}
\label{Complexity.TimeHierarchy.ctrCfg_right}
\leanok
\uses{Complexity.TimeHierarchy.ClockState, Complexity.TimeHierarchy.ctrCfg, Turing.Cfg, Turing.Cfg.workTapeSymbols}
\difficulty{3}
Let $K$ and $W$ be finite Boolean Turing machines, and form a clock configuration from
any clock state, any configuration $c$ of $W$, buffer-tape content and head position,
fixed contents and head positions for the $K$ work tapes, and output word. For every
work-tape index of $W$, the symbol scanned on the corresponding right-block tape of the
clock configuration equals the symbol scanned on that tape in $c$.
\end{lemma}

\begin{lemma}[Counter action on a configuration]
\lean{Complexity.TimeHierarchy.ctrAction_apply}
\label{Complexity.TimeHierarchy.ctrAction_apply}
\leanok
\uses{Complexity.TimeHierarchy.ClockState, Complexity.TimeHierarchy.ctrAction, Complexity.TimeHierarchy.ctrCfg, Turing.Action.apply, Turing.Cfg, Turing.moveInputPos_zero}
\difficulty{5}
Let $K$ and $W$ be finite Boolean Turing machines, let $x$ be a Boolean word, and form a
clock configuration from a $W$-configuration $c$, a buffer tape $\tau$, buffer-head
position $z$, fixed contents and head positions for the $K$ work tapes, clock state
$st$, and output word. Applying a counter action that writes the optional value $w$ at
the buffer head, moves that head by $d$, and sets the state to $q$ produces the clock
configuration with state $q$, buffer tape unchanged if $w$ is absent and otherwise
updated to $w$ at $z$, buffer-head position $z+d$, and all other listed components
unchanged.
\end{lemma}

\begin{definition}[Counter-scan configuration]
\lean{Complexity.TimeHierarchy.cScanCfg}
\label{Complexity.TimeHierarchy.cScanCfg}
\leanok
\uses{Complexity.TimeHierarchy.clockTM, Complexity.TimeHierarchy.ctrCfg, Turing.Cfg, Turing.FinTM.bufferTape}
Let $K$ and $W$ be finite Boolean Turing machines, let $x$ and $s$ be Boolean words, let
$p$ be one of the $|x|+2$ input-head positions, and let $kt$ and $kh$ specify contents
and head positions for the work tapes of $K$. The clock machine configuration has state
$cScan$, input-head position $p$, empty output, and blank work tapes and zero head
positions in the simulated $W$ block. Its $K$ block is given by $kt$ and $kh$; the
intervening buffer tape stores $s$ starting at cell $j-1$, with its head at $j-1$, where
$j\in\mathbb N$.
\end{definition}

\begin{lemma}[Counter scan initialization]
\lean{Complexity.TimeHierarchy.cScanCfg_run}
\label{Complexity.TimeHierarchy.cScanCfg_run}
\leanok
\uses{Complexity.TimeHierarchy.cScanCfg, Complexity.TimeHierarchy.clockTM, Complexity.TimeHierarchy.clockTM_step_some, Complexity.TimeHierarchy.clockTr, Complexity.TimeHierarchy.ctrAction_apply, Complexity.TimeHierarchy.ctrCfg, Complexity.TimeHierarchy.ctrCfg_ctr, Turing.FinTM.bufferTape, Turing.FinTM.bufferTape_left, Turing.FinTM.bufferTape_nat, Turing.MultiTapeTM.runFrom, Turing.MultiTapeTM.runFrom_succ_eq_step, Turing.MultiTapeTM.runFrom_succ_eq_step', Turing.MultiTapeTM.runFrom_zero, Turing.MultiTapeTM.step}
\difficulty{6}
Let $K$ and $W$ be finite Boolean Turing machines, let $x$ and $s$ be Boolean words, and
let $p$ be a valid input-head position for input $x$. Fix functions giving the contents
and head positions of the $K$ work tapes. For every natural number $j$ with $j\le |s|$,
the clock machine, started in its counter-scan state with buffer tape containing $s$ and
buffer-head position $j-1$, reaches after $j+1$ steps the input configuration unchanged,
with the buffer head at position $0$ and the clock state in its input-start state.
\end{lemma}

\begin{lemma}[Rewinding after halting]
\lean{Complexity.TimeHierarchy.kCfg_rewind}
\label{Complexity.TimeHierarchy.kCfg_rewind}
\leanok
\uses{Complexity.TimeHierarchy.cScanCfg, Complexity.TimeHierarchy.cScanCfg_run, Complexity.TimeHierarchy.clockTM, Complexity.TimeHierarchy.clockTM_step_some, Complexity.TimeHierarchy.clockTr, Complexity.TimeHierarchy.ctrAction_apply, Complexity.TimeHierarchy.ctrCfg, Complexity.TimeHierarchy.kCfg, Turing.Cfg, Turing.FinTM.bufferTape, Turing.MultiTapeTM.runFrom, Turing.MultiTapeTM.runFrom_succ_eq_step, Turing.MultiTapeTM.step}
\difficulty{4}
Let $K$ and $W$ be finite Turing machines over the binary alphabet, let $x$ be a binary
input, and let $c$ be a configuration of $K$ on $x$ whose state is absent. Starting the
clock machine from the lifted configuration of $c$ and running it for
$|c.\mathrm{output}|+2$ steps yields a configuration in state $\mathrm{iStart}$, with
input head at the position in $c$, the work tapes of $K$ unchanged, the counter tape
storing $c.\mathrm{output}$ with its head at cell $0$, all work tapes belonging to $W$
blank with their heads at cell $0$, and empty output.
\end{lemma}

\begin{lemma}[Clock machine setup]
\lean{Complexity.TimeHierarchy.clockTM_setup}
\label{Complexity.TimeHierarchy.clockTM_setup}
\leanok
\uses{Complexity.TimeHierarchy.clockTM, Complexity.TimeHierarchy.ctrCfg, Complexity.TimeHierarchy.kCfg_init, Complexity.TimeHierarchy.kCfg_rewind, Complexity.TimeHierarchy.kCfg_run, Turing.FinTM.ComputesInTime, Turing.FinTM.ComputesInTime.output_unique, Turing.FinTM.bufferTape, Turing.FinTM.computesInTime_iff, Turing.FinTM.timed_rewind, Turing.MultiTapeTM.initCfg, Turing.MultiTapeTM.runFrom, Turing.MultiTapeTM.runFrom_add}
\difficulty{5}
Let $K$ and $W$ be finite Boolean Turing machines. If $K$ halts on Boolean input $x$
with output word $s$ after exactly $t_K$ steps, then there are a natural number $a$ and
choices of contents and head positions for the $K$ work tapes such that $a\le
t_K+|s|+|x|+5$, and after $a$ steps the clock machine started on $x$ is in its decrement
state with $W$'s initial configuration on $x$, buffer tape containing $s$ with its head
at position $0$, and empty output.
\end{lemma}
